\documentclass[a4paper,11pt]{article}
\usepackage[utf8]{inputenc}
\usepackage[T1]{fontenc}
\usepackage[left=0.78in,right=0.78in,top=0.65in,bottom=0.65in]{geometry}
\usepackage{lmodern}
\usepackage{microtype}
\usepackage[hidelinks]{hyperref}
\input{glyphtounicode}
\pdfgentounicode=1
\hypersetup{pdftitle={Peng Li -- Application -- FOI-PERS-2026-1039},pdfauthor={Peng Li}}
\pagestyle{empty}
\setlength{\parindent}{0pt}
\setlength{\parskip}{8pt}
\setlength{\emergencystretch}{2em}
\urlstyle{same}
\begin{document}
{\LARGE\textbf{Peng Li}}\\[3pt]
Stockholm, Sweden \quad | \quad +46 79 010 31 96\\
\href{mailto:pengli.2810@gmail.com}{pengli.2810@gmail.com}
\hfill 5 November 2026

Andreas Nöjd and the Recruitment Team\\
Hydroacoustics Unit, Swedish Defence Research Agency (FOI)\\
Stockholm--Kista

\textbf{Application: Forskare inom undervattensakustik med inriktning mot modellering}\\
Reference: FOI-PERS-2026-1039

Dear Andreas Nöjd and Recruitment Team,

I am applying for the researcher position in underwater acoustics with a focus on modelling.
My doctoral research at KTH combines acoustic wave modelling, computational fluid dynamics
(CFD), and fluid--structure interaction (FSI) to understand how unsteady flow and vibrating
structures generate sound. I am scheduled to defend my thesis,
\textit{Computational Aero-acoustics of Snoring with Impact on Obstructive Sleep Apnea Diagnosis},
on \textbf{4 December 2026}. I would like to bring this modelling experience to FOI's work on
underwater sound propagation and acoustic signatures.

My work covers both analytical acoustics and high-fidelity numerical simulation.
In a current manuscript, I developed a resonance model for a bifurcated duct with a compliant
splitter, combining convected plane waves, mode-matched junction impedance, and
frequency-dependent structural compliance. Comparisons with COMSOL finite-element
calculations gave frequency differences within 0.002\% for 27 rigid-splitter resonances and
1.03\% for the first ten soft-splitter common resonances in the tested geometry.
This work has given me practical experience in formulating wave problems, treating boundaries
and interfaces, and checking simplified models against numerical reference solutions.

For the coupled flow and sound-generation problem, I developed three-dimensional
STAR-CCM+ simulations using compressible LES, FEM-based hyperelastic tissue models,
moving overset meshes, and strong two-way FSI. I evaluate acoustic sources using the
Ffowcs Williams--Hawkings analogy and duct Green's functions. Pressure and vorticity SPOD,
spectral analysis, and coherence help me distinguish acoustic resonances from convective
flow structures. This research led to a 2026 \textit{Physics of Fluids} Featured Article.
My experience concerns aeroacoustics; the transferable strengths I would bring to hydroacoustics
are wave modelling, acoustic--structural coupling, source analysis, and numerical verification.

I also bring scientific programming and industrial modelling experience.
At Joyson Safety Systems, I developed Python-based transient initiator models, coupled them
to OpenFOAM simulations of shock-wave propagation, and compared predicted pressures with
measurements. At KTH, I use NAISS HPC resources at PDC and NSC and conduct mesh and
time-step convergence studies. My M.Sc. in Mechanics from TU Darmstadt, completed with
a final grade of 1.33, supports this work with a foundation in fluid mechanics, structural
mechanics, and numerical methods for partial differential equations.

FOI's combination of computational research, measurement, and practical applications is
particularly appealing to me. I would welcome the opportunity to develop tools for shallow-water
propagation and investigate the mechanisms behind underwater acoustic signatures and noise.
I have led analysis and manuscript writing while collaborating with researchers in acoustics,
fluid mechanics, and structural dynamics. I use English for research and scientific writing.
My Swedish is currently \textbf{A2}, and I would be committed to developing the proficiency
required for this role.

Thank you for considering my application. I would welcome a discussion of how my experience
in acoustic modelling, CFD, and FSI could contribute to the Hydroacoustics Unit.

Kind regards,\\[3pt]
\textbf{Peng Li}
\end{document}
